% Part: many-valued-logic
% Chapter: three-valued-logics
% Section: multiple-designation

\documentclass[../../../include/open-logic-section]{subfiles}

\begin{document}

\olfileid{mvl}{thr}{mul}

\olsection{$\True$ మాత్రమే కాక ఇతర విలువలనూ నిర్దేశించడం}

ఇప్పటివరకు చూసిన తర్కాలన్నిటిలో నిర్దేశిత
సత్యమూల్యాల సమితి $V^+ = \{\True\}$;
అంటే ఒక దాని సత్యమూల్యం~$\True$ అయితేనే
దాన్ని నిజంగా పరిగణించాం. కానీ దాని విలువ
~$\False$ కాకపోయినా నిజంగా పరిగణించవచ్చు.
ఉదాహరణకు మాత్రికలో $V^+ = \{\True, \Undef\}$
అని నిర్దేశిస్తే మరో తర్కం లభిస్తుంది.

\begin{defn}
\emph{వైరుధ్యాభాస తర్కం}~$\LogLP$ను కింది
మాత్రికతో నిర్వచిస్తాం:
\begin{enumerate}
  \item చరాల నుంచి $\lnot$, $\land$, $\lor$, $\lif$
  సంయోజకాలతో మాత్రమే నిర్మించిన ప్రామాణిక
  ప్రతిజ్ఞావాక్య భాష $\Lang L_0$ యొక్క ఉపభాష.
  \item సత్యమూల్యాల సమితి $V = \{\True, \Undef, \False\}$.
  \item $\True$, $\Undef$లు నిర్దేశిత విలువలు;
  అంటే $V^+ = \{\True, \Undef\}$.
  \item సత్యమూల్య ప్రమేయాలు బలమైన క్లీని
  తర్కంలో ఉన్నట్లే ఉంటాయి.
\end{enumerate}
\sourcecorrection{OLTEMVLMUL-001}{ఆంగ్ల మూలం
ప్రామాణిక భాషను పేర్కొన్నా, దానిలోని అసత్య
స్థిరాంకానికి విలువ ఇవ్వలేదు. బలమైన క్లీని
మాత్రికతో పంచుకునే నాలుగు సంయోజకాల
చర-ఆధారిత ఉపభాషకే ఇక్కడి నిర్వచనాన్ని
పరిమితం చేశాం.}
\end{defn}

\begin{defn}
హాల్డెన్ యొక్క \emph{అర్థరహిత తర్కం}~$\LogHal$ను
కింది మాత్రికతో నిర్వచిస్తాం:
\begin{enumerate}
  \item చరాల నుంచి $\lnot$, $\land$, $\lor$, $\lif$
  సంయోజకాలతో నిర్మించిన ప్రామాణిక
  ప్రతిజ్ఞావాక్య భాష $\Lang L_0$ ఉపభాషకు
  $1$-స్థానిక సంయోజకం~$+$ను చేర్చిన భాష.
  \item సత్యమూల్యాల సమితి $V = \{\True, \Undef, \False\}$.
  \item $\True$, $\Undef$లు నిర్దేశిత విలువలు;
  అంటే $V^+ = \{\True, \Undef\}$.
  \item సత్యమూల్య ప్రమేయాలు బలహీనమైన క్లీని
  తర్కంలో ఉన్నట్లే ఉంటాయి; అదనంగా
  “అర్థరహితం” అనే సంచాలకం ఉంది:
  \begin{center}
    \begin{tabular}{c|c}
    $\tf{+}$ & \\
    \hline
    $\True$ & $\False$ \\
    $\Undef$ & $\True$ \\
    $\False$ & $\False$
    \end{tabular}
  \end{center}
\end{enumerate}
\sourcecorrection{OLTEMVLMUL-002}{ఆంగ్ల మూలం
ప్రామాణిక భాష అని చెప్పి అసత్య స్థిరాంకాన్ని
జాబితాలోనూ పట్టికల్లోనూ ఇవ్వలేదు. బలహీన
క్లీని తర్కపు నాలుగు సంయోజకాల ఉపభాషకు
అదనపు ఏకస్థానిక సంయోజకాన్ని చేర్చిన
భాషగా ఇక్కడ స్పష్టం చేశాం.}
\end{defn}

సత్య పట్టికలు పంచుకునే క్లీని తర్కాలకు
భిన్నంగా, ఈ తర్కాల్లో \emph{సర్వసత్యాలు ఉన్నాయి}.

\begin{prop}\ollabel{prop:LP-taut-CL}
  సాధారణ నాలుగు-సంయోజక భాషలో $\LogLP$
  సర్వసత్యాలే సాంప్రదాయిక ప్రతిజ్ఞావాక్య
  తర్కపు సర్వసత్యాలు.
\sourcecorrection{OLTEMVLMUL-003}{భాషల
పరిధులు భిన్నంగా ఉండకుండా, ఈ సమానత్వ
వాదాన్ని రెండు తర్కాలకూ ఉమ్మడిగా ఉన్న
నాలుగు-సంయోజక భాగానికి పరిమితం చేశాం.
ప్రామాణిక భాషలోని అసత్య స్థిరాంకం ఈ
మాత్రికలో నిర్వచించబడలేదు.}
\end{prop}

\begin{proof}
  \olref[syn][sub]{prop:mvl-cl} ప్రకారం
  $\Entails[\LogLP] !A$ అయితే
  $\Entails[\LogCL] !A$. తిరుగు దిశను
చూపడానికి, $\pAssign v\colon \PVar \to \{\False, \True,
  \Undef\}$ అనే \tetoken{సత్యమూల్య కేటాయింపులో}{!!a{valuation}}~
  $\pValue v(!A)[\LogKs] = \False$ అయితే,
  $\pAssign {v'}\colon \PVar \to \{\False, \True\}$
  అనే \tetoken{సత్యమూల్య కేటాయింపులో}{!!a{valuation}}~
  $\pValue {v'}(!A)[\LogCL] = \False$ అవుతుందని
  చూపుతాం. దీనితో $\LogLP$కు ఫలితం
స్థాపితమవుతుంది: $\LogKs$, $\LogLP$లకు
ఒకే లక్షణాత్మక సత్యమూల్య ప్రమేయాలు ఉన్నాయి;
$\LogLP$లో నిర్దేశితం కాని ఏకైక సత్యమూల్యం
  $\False$ మాత్రమే. $\LogLP$,~$\LogKs$ల
  మధ్య తేడా ఇదొక్కటే. కనుక
  $\Entails/[\LogLP] !A$ అయితే, ఏదో
  \tetoken{సత్యమూల్య కేటాయింపులో}{!!{valuation}}~
  $\pAssign v$ వద్ద $\pValue v(!A)[\LogLP] = \pValue
  v[\LogKs](!A) = \False$. మనం నిరూపించే
వాదం ప్రకారం $\pValue{v'}[\LogCL](!A) = \False$;
అంటే $\Entails/[\LogCL] !A$.

  ఈ వాదాన్ని స్థాపించడానికి మొదట
  $\pAssign {v'}$ను ఇలా నిర్వచిస్తాం:
  \[
    \pAssign {v'}(p) =
    \begin{cases}
    \True & \text{ఒకవేళ } \pAssign {v}(p) \in \{\True, \Undef\}\\
    \False & \text{లేకపోతే}
  \end{cases}
  \]
  ఇప్పుడు $!A$పై ఆగమన పద్ధతిలో రెండు
విషయాలను చూపుతాం: (a)~$\pValue v(!A)[\LogKs]
  = \False$ అయితే $\pValue {v'}(!A)[\LogCL] = \False$;
  (b)~$\pValue v(!A)[\LogKs] = \True$ అయితే
  $\pValue {v'}(!A)[\LogCL] =
  \True$.
  \begin{enumerate}
    \item ఆగమన ఆధారం: $!A \ident p$.
    \olref[syn][val]{defn:pValue} ప్రకారం
    $\pValue v(!A)[\LogKs] = \pAssign v(p)$,
    అలాగే $\pValue {v'}(!A)[\LogCL] =
    \pAssign {v'}(p)$. మొదటి విలువ
    అసత్యం లేదా సత్యం అయినప్పుడు
    నిర్వచనం ప్రకారం రెండవ విలువ
    వరుసగా అసత్యం లేదా సత్యం;
    అందువల్ల (a), (b) రెండూ నిలుస్తాయి.
    \sourcecorrection{OLTEMVLMUL-004}{ఆంగ్ల
    మూలం చరం అనిర్ణీతంగా ఉన్నప్పుడూ రెండు
    కేటాయింపుల విలువలు సమానమని వ్రాసింది.
    అనిర్ణీత విలువను రెండవ కేటాయింపు
    సత్యంగా మార్చుతుంది; కావలసిన
    అసత్య-సత్య సంరక్షణ వాదాలను మాత్రమే
    ఇక్కడ నిరూపించాం.}

    ఆగమన దశలో కింది సందర్భాలు చూడండి:
    \item $!A \ident \lnot !B$.
    \begin{enumerate}
      \item $\pValue v(\lnot!B)[\LogKs] = \False$
      అనుకుందాం. $\tf{\lnot}[\LogKs]$ నిర్వచనం
      ప్రకారం $\pValue v(!B)[\LogKs]  =
      \True$. ఆగమన పరికల్పనలోని (b) ప్రకారం
      $\pValue {v'}(!B)[\LogCL]  = \True$;
      కాబట్టి $\pValue {v'}(\lnot !B)[\LogCL] = \False$.
      \item $\pValue v(\lnot!B)[\LogKs] = \True$
      అనుకుందాం. $\tf{\lnot}[\LogKs]$ నిర్వచనం
      ప్రకారం $\pValue v(!B)[\LogKs]  =
      \False$. ఆగమన పరికల్పనలోని (a) ప్రకారం
      $\pValue {v'}(!B)[\LogCL]  = \False$;
      కాబట్టి $\pValue {v'}(\lnot !B)[\LogCL] = \True$.
    \end{enumerate}

    \item $!A \ident (!B \land !C)$.
    \begin{enumerate}
      \item $\pValue v(!B \land !C)[\LogKs] = \False$
      అనుకుందాం. $\tf{\land}[\LogKs]$ నిర్వచనం
      ప్రకారం $\pValue v(!B)[\LogKs]  =
      \False$ లేదా $\pValue v(!C)[\LogKs]  =
      \False$. ఆగమన పరికల్పనలోని (a) ప్రకారం
      $\pValue {v'}(!B)[\LogCL]  = \False$
      లేదా $\pValue {v'}(!C)[\LogCL]  = \False$;
      కాబట్టి $\pValue {v'}(!B \land
      !C)[\LogCL] = \False$.
      \item $\pValue v(!B \land !C)[\LogKs] = \True$
      అనుకుందాం. $\tf{\land}[\LogKs]$ నిర్వచనం
      ప్రకారం $\pValue v(!B)[\LogKs]  =
      \True$ మరియు $\pValue v(!C)[\LogKs]  =
      \True$. ఆగమన పరికల్పనలోని (b) ప్రకారం
      $\pValue {v'}(!B)[\LogCL]  = \True$
      మరియు $\pValue {v'}(!C)[\LogCL]  = \True$;
      కాబట్టి $\pValue {v'}(!B \land
      !C)[\LogCL] = \True$.
    \end{enumerate}
    \sourcecorrection{OLTEMVLMUL-005}{ఆంగ్ల మూలం
    సంయోగంలో రెండు సందర్భాల్లోనూ రెండవ
    భాగం స్థానంలో మొదటి భాగాన్నే మళ్లీ
    పేర్కొంది. పట్టిక, ఆగమన పరికల్పన,
    చివరి నిర్ధారణకు సరిపోయేలా రెండవ
    భాగాన్ని స్పష్టం చేశాం.}
  \end{enumerate}

  మిగిలిన రెండు సందర్భాలు ఇదే విధంగా
నిరూపించవచ్చు; వాటిని సాధనలుగా వదిలాం.
మరో విధంగా, పై నిరూపణ $\lnot$,
~$\land$ మాత్రమే కలిగిన అన్ని
\tetoken{సూత్రాలకు}{!!{formula}s} ఫలితాన్ని
స్థాపిస్తుంది. ఇప్పుడు $\LogKs$,
$\LogCL$ రెండింటిలోనూ ఏ $\pAssign v$కైనా
$\pValue v(!B \lor
  !C) = \pValue v(\lnot(\lnot!B \land \lnot !C))$
మరియు $\pValue v(!B \lif
  !C) = \pValue v(\lnot(!B \land \lnot !C))$
అనే వాస్తవాలను వాడవచ్చు.
\end{proof}

\begin{prob}
  \olref[mvl][thr][mul]{prop:LP-taut-CL}
  నిరూపణను పూర్తి చేయండి. అంటే
  $!A \ident (!B \lor !C)$, $!A \ident (!B \lif !C)$
  సందర్భాల్లో (a), (b)లను స్థాపించండి.
\end{prob}

\begin{prob}
ప్రతి సాంప్రదాయిక సర్వసత్యమూ~$\LogHal$లో
సర్వసత్యమని నిరూపించండి.
\end{prob}

సాధారణ నాలుగు-సంయోజక భాషలో సాంప్రదాయిక
తర్కంతో ఒకే సర్వసత్యాలు ఉన్నప్పటికీ, ఈ
తర్కాల పర్యవసాన సంబంధాలు వేరు.
ఉదాహరణకు $\LogLP$ \emph{విస్ఫోటనరహిత}
తర్కం: $\lnot p, p \Entails/ q$.
కాబట్టి విస్ఫోటన సూత్రం
$\lnot !A, !A \Entails !B$ సాధారణంగా
నిలవదు. కొన్ని $!A$, $!B$ సందర్భాల్లో,
ఉదాహరణకు $!B$~సర్వసత్యమైతే, అది నిలుస్తుంది.
\sourcecorrection{OLTEMVLMUL-006}{ఆంగ్ల
మూలంలోని “ఒకే సర్వసత్యాలు” వాదాన్ని
రెండు తర్కాలకూ, సాంప్రదాయిక తర్కానికీ
ఉమ్మడిగా ఉన్న నాలుగు-సంయోజక భాషకే
పరిమితం చేశాం. హాల్డెన్ తర్కానికి
అదనపు సంచాలకం ఉంది కాబట్టి మొత్తం
సూత్ర సమితులను నేరుగా సమానమనలేం.}

\begin{prob}
  కింది సంబంధాల్లో ఏవి (a)~$\LogLP$లో,
  (b)~$\LogHal$లో నిలుస్తాయి?
  ప్రతిదానికీ సత్య పట్టిక ఇవ్వండి.
  \begin{enumerate}
    \item $p, p \lif q \Entails q$
    \item $\lnot q, p \lif q \Entails \lnot p$
    \item $p \lor q, \lnot p \Entails q$
    \item $\lnot p, p \Entails q$
    \item $p \Entails p \lor q$
    \item $p \lif q, q\lif r \Entails p \lif r$
  \end{enumerate}
\end{prob}

$\LogLuk[3]$లో $\Undef$ను నిర్దేశిత
విలువగా తీసుకుంటే ఏమవుతుంది?

\begin{defn}
  \emph{మూడు-విలువల ఆర్-మింగిల్}
  తర్కం~$\LogRM[3]$ను కింది మాత్రికతో
  నిర్వచిస్తాం:
  \begin{enumerate}
    \item $\lfalse$, $\lnot$, $\land$, $\lor$, $\lif$
    సంయోజకాలు గల ప్రామాణిక ప్రతిజ్ఞావాక్య
    భాష $\Lang L_0$.
    \item సత్యమూల్యాల సమితి $V = \{\True, \Undef, \False\}$.
    \item $\True$, $\Undef$లు నిర్దేశిత విలువలు;
    అంటే $V^+ = \{\True, \Undef\}$.
    \item సత్యమూల్య ప్రమేయాలు \L ukasiewicz
    తర్కం~$\LogLuk[3]$లో ఉన్నట్లే ఉంటాయి.
  \end{enumerate}
\end{defn}

\begin{prob}
  కింది సంబంధాల్లో ఏవి $\LogRM[3]$లో నిలుస్తాయి?
  \begin{enumerate}
    \item $p, p \lif q \Entails q$
    \item $p \lor q, \lnot p \Entails q$
    \item $\lnot p, p \Entails q$
    \item $p \Entails p \lor q$
  \end{enumerate}
\end{prob}

నిర్దేశిత విలువలను మార్చడం వల్ల,
వేర్వేరు సత్య పట్టికలు కొన్ని సందర్భాల్లో
ఒకే తర్కాన్ని, అంటే ఒకే అనుగమన సంబంధాన్ని,
ఇవ్వవచ్చు. ఉదాహరణకు గోడెల్ తర్కంలో
~$\{\True\}$ బదులు $V^+ = \{\True,
\Undef\}$ తీసుకుంటే ఇలా జరుగుతుంది.

\begin{prop}\ollabel{prop:gl-udes}
  $V = \{\False, \Undef,\True\}$,
  $V^+=\{\True,
  \Undef\}$, $3$-విలువల గోడెల్ తర్కపు
  సత్యమూల్య ప్రమేయాలు గల మాత్రిక
  సాంప్రదాయిక తర్కాన్ని నిర్వచిస్తుంది.
\end{prop}

\begin{proof}
  సాధన.
\end{proof}

\begin{prob}
  \olref[mvl][thr][mul]{prop:gl-udes}ను
నిరూపించండి. అంటే, గోడెల్ తర్కం మాదిరిగానే
నిర్వచించి, $V^+=\{\True,\Undef\}$గా
మార్చిన తర్కం~$\Log L$లో
$\Gamma \Entails/[\Log L] !B$ అయితే
$\Gamma \Entails/[\LogCL] !B$ అని చూపండి.
\olref[mvl][thr][mul]{prop:LP-taut-CL}
ఆలోచనలను వాడండి; అయితే (a), (b)
లక్షణాలను నిరూపించడానికి బదులుగా
$\pValue v(!A)[\LogGod] = \False$
అయితే, అప్పుడే $\pValue {v'}(!A)[\LogCL] = \False$
అని చూపండి. అందువల్ల
$\pValue
  v(!A)[\LogGod] \in \{\True,\Undef\}$
అయితే, అప్పుడే
$\pValue {v'}(!A)[\LogCL] =
  \True$ అని చూపినట్లే. ఇది ప్రతిపాదనను
ఎందుకు స్థాపిస్తుందో వివరించండి.
\end{prob}

\end{document}
